\worldchapter{Fortschritt und Downtime}{advancement}
\label{ch:fortschritt}

\begin{multicols*}{2}
\section{Grundsatz}
Fortschritt belohnt uberstandene Kosten, Verwundbarkeit und Reputation, nicht stumpfes Wiederholen.

\section{Growth}
Am Ende einer Sitzung erhält ein Charakter 1 Growth, wenn mindestens einer der folgenden Punkte zutraf:
\begin{itemize}[leftmargin=*]
\item sie hat eine groBe Komplikation aus einer Verwundbarkeit akzeptiert
\item sie hat eine \textbf{Critical Consequence} uberlebt, die mit einem Mark verknupft war
\end{itemize}

Bei 3 Growth wahlst du genau eines:
\begin{itemize}[leftmargin=*]
\item eine Fertigkeit von 0 auf 1 erhohen
\item eine Fertigkeit von 1 auf 2 erhohen
\item eine Fertigkeit von 2 auf 3 erhohen
\item eine \textbf{Mark Edge Sentence} erganzen
\item 1 Corruption loschen und 1 neue Obligation erhalten
\end{itemize}

\section{Downtime}
Zwischen groBen Einsätzen erhalt jeder Charakter \textbf{2 Downtime-Actions}.
Wiederholst du dieselbe Aktion in derselben Downtime, bekommt jeder Wiederholungsversuch kumulativ -10.

\section{Downtime-Aktionen}

\section{Downtime-Aktionen}

Während der Downtime kann ein Charakter eine der folgenden Aktionen ausführen:

%\begin{itemize}[leftmargin=*]
%\item \textbf{Bindung pflegen}
%\item \textbf{Wunden versorgen}
%\item \textbf{Kraft schöpfen}
%\item \textbf{Gunst suchen}
%\item \textbf{Nachforschungen anstellen}
%\item \textbf{Arkana sammeln}
%\item \textbf{Reinigung suchen}
%\item \textbf{Ausrüstung beschaffen}
%\end{itemize}

\subsection{Recover}
\[
Will + Healing
\]
Erfolg: 2 Wounds oder 3 Burden oder 1 Wound + 1 Burden + 1 Zustand loschen.\\
Misserfolg: nur 1 Burden loschen.

\subsection{Repair}
\[
Mind oder Deftness + Craft
\]
Erfolg: ein Broken-Gegenstand vollstandig repariert.\\
Starker oder kritischer Erfolg: zwei Gegenstande.\\
Misserfolg: keine Reparatur; erst neues Material, Bezahlung oder Zeit ermoglicht einen weiteren Versuch.

\subsection{Acquire}
\[
Mind oder Appearance + Influence
\]
Erfolg: ein nutzlicher Gegenstand, Dienst oder zeitweiliger Kontakt.\\
Starker oder kritischer Erfolg: zwei Dinge oder ein hochwertiger Vermogensvorteil.\\
Misserfolg: Preis verdoppelt sich, Gefallen wird geschuldet oder Misstrauen wachst.

\subsection{Research}
\[
Mind + Lore oder Arcana
\]
Erfolg: eine klare, handlungsrelevante Wahrheit.\\
Starker oder kritischer Erfolg: zusatzlich eine ausnutzbare Schwache mit +10 einmalig.\\
Misserfolg: falsche Fährte, Verzogerung oder offengelegte Nachforschung.

\subsection{Train}
\[
Will + Resolve
\]
Erfolg: 1 Training Tick speichern.\\
Bei 3 Training Ticks: ein zulassiger Skill-Anstieg wird konkret gemacht, sofern Growth dafur ausgegeben wurde.\\
Misserfolg: kein Tick; bei Uberarbeitung +1 Burden.

\subsection{Work Contact}
\[
Appearance + Influence \text{ oder } Will + Resolve
\]
Erfolg: einen Strained Bond reparieren oder einen konkreten Gefallen sichern.\\
Starker oder kritischer Erfolg: beides.\\
Misserfolg: der Kontakt fordert Preis, Pflicht oder Verzogerung.

\subsection{Ritual Work}
\[
Mind oder Will + Zauberei oder Faith
\]
Erfolg: ein Ritualvorteil, +10 auf einen zukunftigen passenden Cast oder eine Petition.\\
Starker oder kritischer Erfolg: zusatzlich 1 Arcana oder 1 Gunst regenerieren.\\
Misserfolg: instabiles Omen, Spur eines Backlashs oder +1 Burden.

\section{Obligations}
Eine Obligation ist eine dauerhafte Schuld, Pflicht oder ein Dienst.
\begin{itemize}[leftmargin=*]
\item einmal pro Sitzung darf der Guide sie einfordern
\item akzeptieren: +1 Omen und die Komplikation tragen
\item verweigern: +1 Burden und die zugehorige Beziehung verschlechtert sich
\end{itemize}

\section{Auffullen der Ressourcen}
Am Ende der Downtime:
\begin{itemize}[leftmargin=*]
\item Omen auf Max auffullen
\item Arcana Pool und Gunst voll auffullen
\end{itemize}
\end{multicols*}
